% Part: computability
% Chapter: recursive-functions
% Section: normal-form

\documentclass[../../../include/open-logic-section]{subfiles}

\begin{document}

\olfileid{cmp}{rec}{nft}
\olsection{प्रमाण रूप प्रमेय}

\begin{thm}[क्लिनीचे प्रमाण रूप प्रमेय]
\ollabel{thm:kleene-nf}
असा आदिम पुनरावर्ती संबंध $T(e, x, s)$ आणि असे आदिम पुनरावर्ती फलन
$U(s)$ असते की पुढील गुणधर्म पूर्ण होतो: $f$ हे कोणतेही आंशिक पुनरावर्ती
फलन असेल, तर एखाद्या~$e$ साठी सर्व $x$ बाबत
\[
f(x) \simeq U(\umin{s}{T(e, x, s)})
\]
असते.
\end{thm}

\begin{explain}
प्रमाण रूप प्रमेयाची सिद्धता गुंतागुंतीची असली, तरी तिची मूलभूत कल्पना
सोपी आहे. प्रत्येक आंशिक पुनरावर्ती फलनाला एक \emph{निर्देशांक}~$e$
असतो; सहज अर्थाने, तो त्या फलनाचा कार्यक्रम किंवा व्याख्या संकेतबद्ध
करणारी संख्या असते. $f(x) \fdefined$ असेल, तर त्या संगणनाची पद्धतशीर
नोंद करून तिला एखाद्या संख्येने~$s$ संकेतबद्ध करता येते. केवळ $x$ आणि
व्याख्या~$e$ वापरून, $s$ ही संख्या आदान~$x$ वरील~$f$ चे संगणन
संकेतबद्ध करते का, हे आदिम पुनरावर्ती पद्धतीने तपासता येते. म्हणून
``निर्देशांक~$e$ असलेल्या फलनाचे आदान~$x$ साठी संगणन आहे आणि $s$
त्या संगणनाचा संकेतांक आहे'' हा संबंध~$T$ आदिम पुनरावर्ती असतो.
संगणनाची पूर्ण नोंद~$s$ दिली असता, $s$ चा परिणाम म्हणजे~$f(x)$ चे मूल्य;
तेही~$s$ पासून आदिम पुनरावर्ती पद्धतीने मिळवता येते.

प्रमाण रूप प्रमेय दाखवते की कोणत्याही आंशिक पुनरावर्ती फलनाच्या
व्याख्येसाठी एकच अपरिबद्ध शोध पुरेसा असतो. म्हणजे, निर्देशांक~$e$
असलेल्या फलनाचे आदान~$x$ वरील संगणन संकेतबद्ध करणारी संख्या सापडेपर्यंत
आपण सर्व संख्यांमध्ये शोध करू शकतो. आंशिक पुनरावर्ती फलनांची ``नावे''
म्हणून या संख्या~$e$ वापरता येतात आणि प्रमेयातील समीकरणाने परिभाषित
केलेल्या~$f$ फलनासाठी $\cfind{e}$ असे लिहिता येते. एका आंशिक
पुनरावर्ती फलनाला एकापेक्षा अधिक निर्देशांक असू शकतात, हे लक्षात घ्या;
खरे तर प्रत्येक आंशिक पुनरावर्ती फलनाला अनंत निर्देशांक असतात.
\end{explain}
\end{document}

